\chapter{Implementations and Experimental Evaluation}
\label{ch:implementations}
\minitoc

The architectures of Chapter~\ref{ch:architectures} are now brought to
experimental evaluation. Each is implemented on a standard benchmark
and evaluated against an MLP baseline that ignores the structure of the graph,
so that whatever performance separates them is attributable to the relational
inductive bias of Section~\ref{sec:relational-bias}, under which the examples
are related rather than independent. This comparison directly addresses the
objective stated in Section~\ref{sec:objectives-scope}. The chapter first
establishes the methodology common to all three implementations and then
presents each in turn: GCN on the Cora benchmark dataset for node classification,
GAE on Amazon Computers for link prediction, and GIN on NCI1 for graph classification.

\section{Experimental Methodology}
\label{sec:methodology}

All three implementations are built in PyTorch
Geometric~\cite[-2.1pc]{fey2019}, a library that provides graph data structures,
benchmark loaders, and message-passing layers, and all three follow one
experimental pipeline. A benchmark dataset is loaded and prepared;
the GNN and the MLP are instantiated; both are trained under identical optimization
settings, with a validation set monitored to select the best model; and the
selected models are evaluated once on a held-out test set.

The datasets are established benchmarks available through PyTorch Geometric.
For this reason, no exploratory analysis or additional cleaning is performed,
and preparation is limited to the transformations and splits described in the
corresponding section.

No hyperparameter search is performed. Instead, the architectures and optimization
settings are selected to provide stable learning, seeking configurations that avoid
underfitting and overfitting as far as each architecture permits. The goal is to ensure
that neither the GNNs nor the MLPs are disadvantaged by poor configuration,
so that the comparison isolates the effect of incorporating graph structure
through message passing.

Reproducibility is enforced through controlled execution conditions.
All experiments are executed on CPU, avoiding hardware-dependent
nondeterminism introduced by parallel GPU operations. Under this setting,
every stochastic component of a run --- the initialization of parameters,
the sampling of splits, the drawing of negative examples, and other
operations that consume randomness --- is governed by a single seed,
so that a run is reproducible once its seed is fixed.
Each experiment is repeated over a fixed set of ten seeds. Reporting
performance over multiple seeds provides a more robust assessment
than a single run by accounting for the variability induced
by stochastic training. For each model, the evaluation metric is summarized by
its median, mean, and standard deviation across the ten runs. The figures show
the run corresponding to the median performance, ensuring that every figure
corresponds to an actual run rather than to an aggregate that
does not correspond to any individual run.

The companion repository~\cite{macias2026} contains the implementation required
to reproduce this chapter: the configurations, the dataset loading
and preparation, the training and evaluation pipeline, the scripts that generate every
figure and compute every metric, and the exact dependency manifest. 

\section{Node Classification with GCN}
\label{sec:gcn-implementation}

The first implementation addresses node classification. GCN
(Section~\ref{sec:gcn}) and an MLP are trained to predict the topic of
each publication in the Cora citation network.

\subsection{The Cora Dataset}
\label{sec:cora-dataset}

Cora is a citation network of machine-learning publications~\cite{sen2008}.
Nodes represent papers, each edge represents an undirected citation
between two papers, each node carries a binary bag-of-words vector
indicating which words from a fixed dictionary occur in the paper, and each node
is labeled with one of seven topics. Although citations are inherently directed,
the benchmark represents them as undirected edges to allow the information
to be propagated in both directions during message passing.
Table~\ref{tab:cora-stats} summarizes the dataset.

\begin{table}
  \caption[Cora dataset.]{%
    Summary of the Cora benchmark dataset.
  }
  \label{tab:cora-stats}
  \begin{center}
    \small
    \begin{tabular}{lr}
      \toprule
      Property & Value \\
      \midrule
      Nodes                        & 2{,}708 \\
      Edges                        & 5{,}278 \\
      Average degree               & 3.9 \\
      Node feature dimension       & 1{,}433 \\
      Classes                      & 7 \\
      Split (train/val/test nodes) & 140 / 500 / 1{,}000 \\
      \bottomrule
    \end{tabular}
  \end{center}
\end{table}

The dataset is provided as a benchmark through PyTorch Geometric, so no
additional cleaning is performed. The only transformation applied to the node
features is a row-wise L1 normalization.

The predefined data split is adopted without modification. The dataset ships with
a partition of 140 training nodes (balanced with twenty per class), 500 validation nodes,
and 1{,}000 test nodes. The training partition therefore contains only about five percent
of the nodes in the graph. The remaining 1{,}068 nodes belong to none of the three
partitions but remain part of the graph, contributing structural information during
message passing without belonging to any supervised split.
This corresponds to the transductive setting described in
Section~\ref{sec:transductive-inductive}, where training has access to the
full graph structure and all node features, while supervision is provided
only through the labeled training nodes.

\subsection{Model and Training}
\label{sec:gcn-training}

Table~\ref{tab:gcn-model} specifies both models. The GCN stacks two
layers of~\eqref{eq:gcn-matrix}, the first serving as the hidden layer and the second
as the output layer, giving every node a two-hop receptive field
(Section~\ref{sec:receptive-field}). The MLP consists of a single hidden layer of
width 32 followed by the output layer, resulting in approximately twice as many
trainable parameters as the GCN --- wider MLP configurations
overfit the 140 training labels more severely.

\begin{table}
  \caption[Model configuration, node classification.]{%
    Model configuration for node classification on Cora.
  }
  \label{tab:gcn-model}
  \begin{center}
    \small
    \begin{tabular}{lcc}
      \toprule
      Component            & GCN               & MLP \\
      \midrule
      Layers               & 2 GCN layers & 2 MLP layers \\
      Hidden dimension     & 16                & 32 \\
      Hidden activation        &\multicolumn{2}{c}{ReLU} \\      
      Output dimension       &\multicolumn{2}{c}{7} \\
      Trainable parameters & 23{,}063          & 46{,}119 \\
      \bottomrule
    \end{tabular}
  \end{center}
\end{table}

Table~\ref{tab:gcn-training} specifies the optimization settings, identical for
both models. Both produce class probabilities through a softmax transformation
and are trained using cross-entropy loss and the same
regularization, so the comparison isolates the effect of incorporating graph
structure through message passing. Dropout~\cite[-2pc]{srivastava2014} and
L2 regularization (weight decay) are used for regularization.
Training is performed in full-batch mode, where the entire graph is processed
at every optimization step, as permitted by the transductive setting.

\begin{table}
  \caption[Training configuration, node classification.]{%
    Training configuration for node classification on Cora, identical for
    both models.
  }
  \label{tab:gcn-training}
  \begin{center}
    \small
    \begin{tabular}{lr}
      \toprule
      Parameter     & Value \\
      \midrule
      Loss             & Cross-entropy \\
      Optimizer     & Adam \\
      Learning rate & 0.01 \\
      Weight decay  & $5 \times 10^{-4}$ \\
      Dropout       & 0.5 \\
      Epochs        & 200 \\
      Batch size    & full graph \\
      \bottomrule
    \end{tabular}
  \end{center}
\end{table}

The validation set is evaluated at every epoch, and the parameters achieving
the highest validation accuracy are retained. The results reported in the
next subsection are obtained from the selected checkpoint evaluated on the test nodes.

\subsection{Results}
\label{sec:gcn-results}

Figure~\ref{fig:node-loss} shows the loss over training. Both models reduce
the training loss, but the MLP reduces it faster and to a lower value.
On the validation set the two diverge. The GCN validation loss continues its
decreasing trend throughout training, whereas the MLP validation
loss settles early while its training loss continues to decrease,
suggesting overfitting. By incorporating information from neighboring nodes
through message passing, the GCN exhibits substantially less separation
between training and validation loss.

Performance is measured by accuracy over the held-out nodes. Accuracy provides
a global measure of classification effectiveness, while the learned embeddings
are examined separately to assess whether the representations organize nodes
according to their classes. Table~\ref{tab:gcn-results} reports the results.

\newpage

\begin{figure}
  \centering
  \includegraphics{node-loss}
  \caption[Training and validation loss on Cora.]{%
    Training (left) and validation (right) loss per epoch on Cora for GCN and
    the MLP. Each panel is scaled to its own data.
  }
  \label{fig:node-loss}
\end{figure}

\begin{table}
  \caption[GCN against the MLP on Cora.]{%
    Test accuracy on Cora, over ten seeds.
  }
  \label{tab:gcn-results}
  \begin{center}
    \small
    \begin{tabular}{lccc}
      \toprule
      Model & Median & Mean & Std \\
      \midrule
      GCN & 0.8170 & 0.8154 & 0.0061 \\
      MLP & 0.5920 & 0.5876 & 0.0119 \\
      \bottomrule
    \end{tabular}
  \end{center}
\end{table}

The GCN reaches a median accuracy of 0.8170, compared with 0.5920 for
the MLP, an improvement of approximately 22.5 percentage points despite
using about half as many trainable parameters. The MLP relies exclusively on
the node features, whereas the GCN enriches those representations through
message passing. This result is consistent with the smoothing mechanism of
GCNs described in Section~\ref{sec:gcn}, which is effective when
neighboring nodes tend to share the same class.

Figure~\ref{fig:node-embeddings} illustrates this difference. It projects the learned
test-node representations to two dimensions using UMAP~\cite{mcinnes2018}.
The GCN representations exhibit a clearer clustering by class, whereas the MLP
representations remain more diffuse with greater overlap between classes.
Because UMAP preserves local neighborhoods rather than global distances or
cluster sizes, only the local clustering within each panel is informative.

\begin{figure}
  \centering
  \includegraphics[width=3.2in]{node-embeddings}
  \caption[Learned representations on Cora.]{%
    Two-dimensional UMAP projection of the test-node representations learned
    on Cora by GCN (left) and by the MLP (right), colored by class.
    Each panel is an independent projection; only the grouping within a panel
    is informative.
  }
  \label{fig:node-embeddings}
\end{figure}

\FloatBarrier

\section{Link Prediction with GAE}
\label{sec:gae-implementation}

The second implementation addresses link prediction. GAE
(Section~\ref{sec:gae}) and an MLP-based encoder with the same decoder
are trained to predict co-purchase links on the Amazon Computers graph.

\subsection{The Amazon Computers Dataset}
\label{sec:amazon-dataset}

Amazon Computers is a segment of the Amazon co-purchase
graph~\cite{shchur2018}. Each node is a product, each edge is undirected
and connects two products frequently bought together, and each node carries
a bag-of-words vector extracted from product reviews. The task uses no node labels: the
supervision is the edge set itself. Table~\ref{tab:amazon-stats} summarizes
the dataset.

\begin{table}
  \caption[Amazon Computers dataset.]{%
    Summary of the Amazon Computers benchmark dataset.
  }
  \label{tab:amazon-stats}
  \begin{center}
    \small
    \begin{tabular}{lr}
      \toprule
      Property & Value \\
      \midrule
      Nodes                        & 13{,}752 \\
      Edges                        & 245{,}861 \\
      Average degree               & 35.76 \\
      Node feature dimension       & 767 \\
      Split (train/val/test edges) & 85\% / 5\% / 10\% \\
      \bottomrule
    \end{tabular}
  \end{center}
\end{table}

The dataset is provided as a benchmark through PyTorch Geometric, so no
additional cleaning is performed. The only transformation applied to the node
features is a row-wise L1 normalization.

The split is defined over edges. A seeded sampling assigns 85\% of the edges
to training, 5\% to validation, and 10\% to test; validation and test edges
are removed from the graph on which the encoder operates, so message passing
runs on the training structure only. Both the validation and test sets are augmented
with an equal number of sampled non-edges, the negative examples introduced in
Section~\ref{sec:encoder-decoder}, so that evaluation is a binary
classification over held-out pairs, with the non-edges drawn to be genuine ---
absent from the full graph and shared by no two evaluation sets. At this scale,
the evaluation sets contain tens of thousands of pairs each, which keeps the
variability of the resulting estimates low. The training setting remains transductive:
all nodes and their features are visible throughout, while only the evaluation edges
are held out from the training graph.

\subsection{Model and Training}
\label{sec:gae-training}

Table~\ref{tab:gae-model} specifies both models. The GAE encoder stacks two
GCN layers, one as the hidden layer and the other as the output embedding layer;
the MLP encoder consists of two hidden layers applied to each node
independently, followed by an output embedding layer, resulting in about twice
as many parameters as the GAE encoder. Both feed the same inner product
decoder~\eqref{eq:inner-product-decoder}, which has no parameters of its own,
so the encoder is what differs between the two models.

\begin{table}
  \caption[Model configuration, link prediction.]{%
    Model configuration for link prediction on Amazon Computers.
  }
  \label{tab:gae-model}
  \begin{center}
    \small
    \begin{tabular}{lcc}
      \toprule
      Component            & GAE          & MLP \\
      \midrule
      Encoder              & 2 GCN layers & 3 MLP layers \\
      Hidden dimensions    & 64           & 128 / 64 \\
      Hidden activation       &\multicolumn{2}{c}{ReLU} \\
      Embedding dimension  &\multicolumn{2}{c}{32} \\
      Decoder              & \multicolumn{2}{c}{inner product} \\
      Trainable parameters   & 51{,}232     & 108{,}640 \\
      \bottomrule
    \end{tabular}
  \end{center}
\end{table}

Table~\ref{tab:gae-training} specifies the optimization, identical for both
models. Negative examples are resampled at every epoch, with one non-edge
sampled per training edge. The loss is the binary cross-entropy between the
reconstructed probabilities~\eqref{eq:gae-reconstruction} and the observed
adjacency over the sampled pairs, with no regularization.
The choice of 400 epochs is deliberate: at that point both models are
still improving slightly without clear signs of overfitting, as the loss curves
in the next subsection confirm.

\begin{table}
  \caption[Training configuration, link prediction.]{%
    Training configuration for link prediction on Amazon Computers, identical
    for both models.
  }
  \label{tab:gae-training}
  \begin{center}
    \small
    \begin{tabular}{lr}
      \toprule
      Parameter         & Value \\
      \midrule
      Loss             & Binary cross-entropy \\
      Optimizer         & Adam \\
      Learning rate     & 0.01 \\
      Weight decay      & 0 \\
      Dropout           & 0 \\
      Epochs            & 400 \\
      Batch size        & full graph \\
      Negative sampling & 1:1, resampled per epoch \\
      \bottomrule
    \end{tabular}
  \end{center}
\end{table}

The validation set is evaluated at every epoch, and the parameters achieving
the highest validation AUC are retained. The results reported in the
next subsection are obtained from the selected checkpoint evaluated on the test pairs.

\subsection{Results}
\label{sec:gae-results}

Figure~\ref{fig:link-loss} shows the reconstruction loss over the 400 epochs
of training. Both models improve over training, on both the training and
validation sets, and neither shows signs of overfitting: the validation loss
follows the same decreasing trend as the training loss until the last epoch.
The GAE loss remains below the MLP loss from early in training and stays there.
Among the three tasks, this is the one in which the MLP is a genuine competitor:
it learns the problem, so the comparison is between two models that both
perform effectively.

\begin{figure}
  \centering
  \includegraphics{link-loss}
  \caption[Training and validation loss on Amazon Computers.]{%
    Training (left) and validation (right) loss per epoch on Amazon Computers
    for GAE and the MLP-based encoder. Each panel is scaled to its own data.
  }
  \label{fig:link-loss}
\end{figure}

Two threshold-free metrics summarize performance over the held-out pairs.
The AUC is the probability that a random true edge is scored above a random
non-edge; the average precision (AP) summarizes the precision--recall curve,
giving more importance to retrieving true edges among the highest-scored
pairs. Table~\ref{tab:gae-results} reports them.

\begin{table}
  \caption[GAE against the MLP-based encoder on Amazon Computers.]{%
    Link prediction on Amazon Computers, over ten seeds.
  }
  \label{tab:gae-results}
  \begin{center}
    \small
    \begin{tabular}{llccc}
      \toprule
      Metric & Model & Median & Mean & Std \\
      \midrule
      AUC & GAE      & 0.9389 & 0.9352 & 0.0077 \\
          & MLP & 0.8778 & 0.8567 & 0.0297 \\
      \midrule
      AP  & GAE      & 0.9372 & 0.9346 & 0.0071 \\
          & MLP & 0.8741 & 0.8589 & 0.0248 \\
      \bottomrule
    \end{tabular}
  \end{center}
\end{table}

Figure~\ref{fig:link-roc} shows the two ROC curves on the test pairs.
The GAE curve lies above the MLP curve,
with the largest separation at low false-positive rates. In this regime,
the GAE achieves higher true-positive rates for the same false-positive rate,
indicating better discrimination between true and false links.

\begin{figure}
  \centering
  \includegraphics[width=2.2in]{link-roc}
  \caption[ROC curves on Amazon Computers.]{%
    ROC curves on the Amazon Computers test pairs for GAE and the MLP-based
    encoder, with each AUC annotated; the dashed diagonal is
    the performance of a random score.
  }
  \label{fig:link-roc}
\end{figure}

Although link prediction scores are often used to identify the most likely candidate links,
enumerating the top co-purchases would add nothing to the comparison pursued here.

The GAE reaches a median AUC of 0.9389 and a median AP of 0.9372, compared with
0.8778 and 0.8741 for the MLP, an advantage of approximately 6 percentage points on
both metrics, with roughly half the parameters. The margin is the smallest of
the three tasks, and the reason is instructive: the product features alone
provide substantial predictive signal for co-purchase, so an encoder that ignores
the graph already ranks pairs well. Message passing improves the ranking further by
encoding structural information from neighboring nodes into the embeddings.

\section{Graph Classification with GIN}
\label{sec:gin-implementation}

The third implementation addresses graph classification. GIN
(Section~\ref{sec:gin}) and an MLP are trained to predict whether
a chemical compound is active or inactive in the NCI1 benchmark.

\subsection{The NCI1 Dataset}
\label{sec:nci1-dataset}

NCI1 is a benchmark of chemical compounds from the National Cancer
Institute that were screened for activity against non-small cell lung
cancer~\cite{wale2008}. Each graph is a compound: its nodes are atoms, each
carrying a one-hot vector that encodes the atom type; each edge is undirected
and represents a chemical bond between two atoms; and each graph is labeled
active or inactive, with the two classes in near-equal proportion.
Table~\ref{tab:nci1-stats} summarizes the dataset.

\begin{table}
  \caption[NCI1 dataset.]{%
    The NCI1 dataset as served by PyTorch Geometric.
  }
  \label{tab:nci1-stats}
  \begin{center}
    \small
    \begin{tabular}{lr}
      \toprule
      Property & Value \\
      \midrule
      Graphs                        & 4{,}110 \\
      Average nodes per graph       & 29.87 \\
      Average edges per graph       & 32.3 \\
      Average degree                & 2.16 \\
      Node feature dimension        & 37 \\
      Classes                       & 2 \\
      Class balance                 & 2{,}053 / 2{,}057 \\
      Split (train/val/test graphs) & 80\% / 10\% / 10\% \\
      \bottomrule
    \end{tabular}
  \end{center}
\end{table}

The dataset is provided as a benchmark through PyTorch Geometric, so no
additional cleaning is performed. No additional transformation is required,
since node features are already one-hot encoded.

The split is defined over graphs. A seeded, stratified sampling assigns 80\%
of the graphs to training, 10\% to validation, and 10\% to test, preserving
the class proportions in each subset. This allocation balances the need to maximize
the number of graphs available for parameter learning with the need to retain
evaluation subsets large enough to reduce the variability of the resulting estimates.
The setting is inductive (Section~\ref{sec:transductive-inductive}): the graphs on
which the model is evaluated are disjoint from those on which it is trained,
and the shared parameters of the message-passing layers are what allow
the model to generalize to graphs it has never seen.

\subsection{Model and Training}
\label{sec:gin-training}

Table~\ref{tab:gin-model} specifies both models. GIN stacks three layers of the
update~\eqref{eq:gin}, each with a two-layer MLP and a fixed $\epsilon = 0$ as its
update function; the MLP baseline applies three layers to every node independently,
which gives it about twice as many trainable parameters as GIN. Each update function
applies batch normalization~\cite{ioffe2015} before the activation to stabilize the
intermediate representations during training; the MLP baseline applies the same
normalization between its layers, so this component is also shared by both models.
Both models are then followed by the shared sum readout~\eqref{eq:readout} and the
same classification head, making the graph-level pipeline identical
except for the message passing.

\begin{table}
  \caption[Model configuration, graph classification.]{%
    Model configuration for graph classification on NCI1.
  }
  \label{tab:gin-model}
  \begin{center}
    \small
    \begin{tabular}{lcc}
      \toprule
      Component            & GIN               & MLP \\
      \midrule
      Layers                    & 3 GIN layers & 3 MLP layers \\
      Hidden dimensions     & 64                & 256 / 128 / 64 \\
      $\epsilon^{(k)}$     & 0 (fixed)         & --- \\
      Normalization        & \multicolumn{2}{c}{batch normalization} \\
      Hidden activation      & \multicolumn{2}{c}{ReLU} \\
      Readout              & \multicolumn{2}{c}{sum} \\
      Classification head  & \multicolumn{2}{c}{1 linear} \\
      Output dimension   & \multicolumn{2}{c}{2} \\
      Trainable parameters & 24{,}130          & 51{,}906 \\
      \bottomrule
    \end{tabular}
  \end{center}
\end{table}

Table~\ref{tab:gin-training} specifies the optimization,
identical for both models. Both models produce class probabilities
through a softmax transformation and are trained using cross-entropy loss
with no regularization. Training uses a mini-batch size of 32, since the dataset
consists of many independent graphs rather than a single graph.

\begin{table}
  \caption[Training configuration, graph classification.]{%
    Training configuration for graph classification on NCI1, identical for
    both models.
  }
  \label{tab:gin-training}
  \begin{center}
    \small
    \begin{tabular}{lr}
      \toprule
      Parameter     & Value \\
      \midrule
      Loss             & Cross-entropy \\
      Optimizer     & Adam \\
      Learning rate & 0.01 \\
      Weight decay  & 0 \\
      Dropout       & 0 \\
      Epochs        & 200 \\
      Batch size    & 32 \\
      \bottomrule
    \end{tabular}
  \end{center}
\end{table}

The validation set is evaluated at every epoch, and the parameters achieving
the highest validation accuracy are retained. The results reported in the
next subsection are obtained from the selected checkpoint evaluated
on the test graphs.

\subsection{Results}
\label{sec:gin-results}

Figure~\ref{fig:graph-loss} shows the loss over training. The training dynamics
differ substantially between the two models. The GIN training loss decreases steadily
throughout training. Its validation loss reaches a minimum after roughly 50 epochs and
then increases gradually while becoming more variable. This rise does not by itself indicate
a loss of classification performance: the cross-entropy grows when
the model becomes more confident on the graphs it misclassifies.
Since model selection is based on validation accuracy rather than on validation loss,
training is not stopped at that minimum, and the retained checkpoint is the one that classifies
the validation graphs best. In contrast, the MLP training loss decreases for only
a few epochs before flattening near 0.62, where it remains for the rest of training;
its validation loss exhibits the same behavior. The absence of a gap between training
and validation loss indicates that the MLP does not overfit; instead, its early plateau
is consistent with underfitting, suggesting that the available node features alone
are insufficient for further improvement under this architecture.

\begin{figure}
  \centering
  \includegraphics{graph-loss}
  \caption[Training and validation loss on NCI1.]{%
    Training (left) and validation (right) loss per epoch on NCI1 for GIN and
    the MLP. Each panel is scaled to its own data.
  }
  \label{fig:graph-loss}
\end{figure}

Performance is measured by accuracy over the held-out graphs. Accuracy provides
a global measure of classification effectiveness, while the confusion matrices
complement this metric by showing how errors are distributed across classes.
This reveals potential differences in classification behavior between active and
inactive compounds. Table~\ref{tab:gin-results} reports the test accuracy across
ten seeds.

\begin{table}
  \caption[GIN against the MLP on NCI1.]{%
    Test accuracy on NCI1, over ten seeds.
  }
  \label{tab:gin-results}
  \begin{center}
    \small
    \begin{tabular}{lccc}
      \toprule
      Model    & Median & Mean & Std \\
      \midrule
      GIN      & 0.8199 & 0.8168 & 0.0166 \\
      MLP     & 0.6496 & 0.6477 & 0.0277 \\
      \bottomrule
    \end{tabular}
  \end{center}
\end{table}

GIN reaches a median accuracy of 0.8199, compared with 0.6496 for the MLP,
an advantage of approximately 17 percentage points despite using roughly half
the parameters. The gap reflects the information available to each
architecture. With one-hot atom features and no access to connectivity before
the readout, the structure-blind MLP can distinguish graphs only through their
node attributes; graphs with the same multiset of node features receive identical
representations regardless of their connectivity. GIN, in contrast, aggregates
neighborhood information before the graph-level readout, allowing its
representations to incorporate structural patterns.
The results show that message passing provides an advantage when
there is an informative underlying graph structure.

\begin{figure}
  \centering
  \includegraphics{graph-confusion}
  \caption[Confusion matrices on NCI1.]{%
    Row-normalized confusion matrices on the NCI1 test graphs for GIN (left)
    and the MLP (right), under a shared color scale.
  }
  \label{fig:graph-confusion}
\end{figure}

Figure~\ref{fig:graph-confusion} decomposes the errors. GIN classifies both
active and inactive compounds correctly in approximately four out of five
cases. The MLP achieves substantially better performance on inactive compounds
than on active ones, for which its predictions are close to chance. This behavior is
consistent with the limitations of a structure-blind model when the task
requires information beyond node attributes.